\begin{theorem}[Smith 不变量的有限 \texorpdfstring{$p$}{p}-窗完全恢复与顶层行列式指数]\label{thm:xi-smith-finite-pwindow-complete-recovery-determinant-index}
设 \(A\in M_{m\times n}(\ZZ)\) 满足 \(\operatorname{rank}_{\QQ}(A)=m\)，其 Smith 不变量为
\[
d_1\mid d_2\mid\cdots\mid d_m.
\]
固定素数 \(p\)。对 \(k\ge 1\) 定义
\[
K_p(k):=\#\ker(A\bmod p^k),
\qquad
f_p(k):=\log_p K_p(k)-k(n-m),
\]
并记
\[
e_i(p):=v_p(d_i),
\qquad
E_p:=\max_{1\le i\le m}e_i(p).
\]
则成立：
\begin{enumerate}
  \item 有限窗口 \(\bigl(K_p(k)\bigr)_{1\le k\le E_p}\) 已充要决定 \(p\)-初级 Smith 多重集 \(\{e_i(p)\}_{i=1}^{m}\)；对一切 \(k>E_p\) 的数据均为冗余。
  \item 极限
  \[
  \nu_p(A):=\lim_{k\to\infty}f_p(k)
  \]
  存在，且满足
  \[
  \nu_p(A)
  =
  \sum_{i=1}^{m}e_i(p)
  =
  v_p\!\Bigl(\prod_{i=1}^{m}d_i\Bigr).
  \]
  \item 若 \(A\) 为方阵且 \(\det A\neq 0\)，则
  \[
  \nu_p(A)=v_p(\det A).
  \]
\end{enumerate}
\end{theorem}
\begin{proof}
由定理 \ref{thm:killo-smith-loss-spectrum}，
\[
f_p(k)=\sum_{i=1}^{m}\min\bigl(k,e_i(p)\bigr)
\qquad(k\ge 1),
\]
并且
\[
\Delta_p(k):=f_p(k)-f_p(k-1)=\#\{i:\ e_i(p)\ge k\}.
\]
若 \(k>E_p\)，则对所有 \(i\) 都有 \(e_i(p)<k\)，故 \(\Delta_p(k)=0\)。因此 \(f_p\) 在 \(k\ge E_p\) 后恒定，全部 \(k>E_p\) 的数据均不再携带新的信息。另一方面，由
\[
K_p(k)=p^{\,k(n-m)+f_p(k)}
\]
可由窗口 \(\bigl(K_p(k)\bigr)_{1\le k\le E_p}\) 唯一恢复 \(\bigl(f_p(k)\bigr)_{1\le k\le E_p}\)，进而恢复
\[
\Delta_p(k)=f_p(k)-f_p(k-1)
\qquad(1\le k\le E_p),
\qquad
\Delta_p(k)=0
\quad(k>E_p).
\]
于是对每个 \(\ell\ge 1\)，都有
\[
\#\{i:\ e_i(p)=\ell\}
=
\Delta_p(\ell)-\Delta_p(\ell+1),
\]
故整个多重集 \(\{e_i(p)\}_{i=1}^{m}\) 被唯一恢复。

反之，若仅知某个长度 \(L<E_p\) 的初始窗口，则无法唯一确定该多重集。事实上，取某个满足 \(e_j(p)=E_p\) 的指标 \(j\)，把它替换为 \(e_j'(p):=E_p+1\)，其余赋值保持不变。则对所有 \(k\le L\) 都有
\[
\min(k,E_p)=\min(k,E_p+1)=k,
\]
故修改前后的 \(f_p(k)\) 与 \(K_p(k)\) 完全一致，但对应的 \(p\)-初级 Smith 多重集不同。于是长度短于 \(E_p\) 的窗口不可能给出无损恢复，第 \(1\) 条成立。

第 \(2\) 条由平台稳定性立即得到：当 \(k\ge E_p\) 时，
\[
f_p(k)=\sum_{i=1}^{m}e_i(p),
\]
故极限存在并等于该和值。又因 \(\prod_i d_i\) 的 \(p\)-赋值正是 \(\sum_i e_i(p)\)，故
\[
\nu_p(A)=v_p\!\Bigl(\prod_{i=1}^{m}d_i\Bigr).
\]

若 \(A\) 为方阵且 \(\det A\neq 0\)，则 Smith 乘积满足
\[
\prod_{i=1}^{m}d_i=|\det A|,
\]
故
\[
\nu_p(A)=v_p(|\det A|)=v_p(\det A).
\]
第 \(3\) 条成立。证毕。
\end{proof}
\leanverified{paper\_xi\_smith\_finite\_pwindow\_complete\_recovery\_determinant\_index}

\endinput
